\chapter{概念设计}
% 各模块概念设计按照：机械、电控、视觉三部分组织，便于后续详细设计和验证逻辑衔接。

\section{机械概念设计}

在RoboMaster比赛中，空中机器人的作用日益凸显。得益于赛季规则与地形的演变，无人机在赛场上的战术价值持续提升：能够快速击打前哨站、击杀敌方地面关键单位、攻打敌方基地，凭借高机动性在场上灵活执行各类任务。为充分发挥空中机器人的战术潜力，对其机械系统提出以下核心需求：机体需满足检录收纳尺寸限制并实现快速部署，具备充沛的续航与稳定的飞行能力，搭载精确的发射系统与大角度云台以覆盖广阔的打击范围。据此，我们为无人机机械系统制定了整体功能设计指标，如\tabref{tab:无人机整体功能设计}所示。

\begin{longtable}[htbp!]{ m{4cm}<{\centering}  m{9cm}<{\centering} }
\hline
种类  &  指标 \\
\hline
\endfirsthead
\hline
种类  &  指标 \\
\hline
\endhead
\hline
\multicolumn{2}{r}{接下页}
\endfoot
\caption{无人机整体功能设计}
\label{tab:无人机整体功能设计}
\endlastfoot
 \hline
 可折叠收纳 & 机臂与桨保可折叠，折叠后尺寸满足最大收纳尺寸检录要求，便于运输与快速部署 \\
 \hline
 稳定飞行能力 & 动力系统性能强劲，抗干扰能力优秀，飞行平稳可控 \\
 \hline
 定点定位系统 & 搭载GPS、Livox MID-360雷达与激光检测模块，感知自身位姿，保持定点稳定飞行 \\
 \hline
 充沛续航时间 & 实战场景下飞行时间达到7min以上 \\
 \hline
 轻量化高强度机身 & 碳纤维结构，整体刚性强且质量轻 \\
 \hline
 精确发射系统 & 高频顺畅发射不卡弹，10m距离散度控制在小装甲板范围内 \\
 \hline
 大角度双轴云台 & Yaw轴水平$180^{\circ}$旋转，Pitch轴俯仰角满足近距离开符需求，给云台手提供宽广视野 \\
 \hline
\end{longtable}
